%------------------------------------------------------------------------------%
\begin{question}
Escalone a matriz
\(
  A=
  \begin{bmatrix}
    1  & 2  & -1 \\
    2  & 5  & 1  \\
    -1 & -1 & 2
  \end{bmatrix}
\)
\end{question}

%------------------------------------------------------------------------------%
\begin{resolution}{Solução}
  Utilizamos operações elementares sobre as linhas\\
  até obter uma matriz em forma escalonada
\end{resolution}

%------------------------------------------------------------------------------%
\begin{resolution}{Primeira Coluna}
\begin{columns}
\column{0.3\linewidth}
  \[
  \begin{bmatrix}
    1  & 2  & -1 \\
    2  & 5  & 1  \\
    -1 & -1 & 2
  \end{bmatrix}
  \]
  \pause
\column{0.7\linewidth}
  Eliminamos os elementos abaixo do primeiro pivô
  \\ \pause
  \(L_2 \leftarrow L_2-2 L_1\)
  \\ \pause
  \(L_3\leftarrow L_3+L_1\)
\end{columns}
\vspace{-\baselineskip}
\[
  \pause
  L'_2
  \;=\;
  L_2 - 2 L_1
  \pause
  \;=\;
  \begin{bmatrix}
    2  & 5  & 1
  \end{bmatrix}
  - 2
  \begin{bmatrix}
    1  & 2  & -1
  \end{bmatrix}
  \pause
  \;=\;
  \begin{bmatrix}
    0 & 1 & 3
  \end{bmatrix}
\]
\[
  \pause
  L'_3
  \;=\;
  L_3 + L_1
  \pause
  \;=\;
  \begin{bmatrix}
    -1 & -1 & 2
  \end{bmatrix}
  +
  \begin{bmatrix}
    1  & 2  & -1
  \end{bmatrix}
  \pause
  \;=\;
  \begin{bmatrix}
    0 & 1 & 1
  \end{bmatrix}
\]
\pause
\[
  \begin{bmatrix}
    1 & 2 & -1 \\
    0 & 1 & 3  \\
    0 & 1 & 1
  \end{bmatrix}
\]
\end{resolution}

%------------------------------------------------------------------------------%
\begin{resolution}{Segunda Coluna}
\begin{columns}
\column{0.3\linewidth}
  \[
    \begin{bmatrix}
      1 & 2 & -1 \\
      0 & 1 & 3  \\
      0 & 1 & 1
    \end{bmatrix}
  \]
  \pause
\column{0.7\linewidth}
  Eliminamos os elementos abaixo do segundo pivô
  \\ \pause
  \(L_3\leftarrow L_3-L_2\)
\end{columns}
\vspace{-\baselineskip}
\[
  \pause
  L'_3
  \;=\;
  L_3 - L_2
  \pause
  \;=\;
  \begin{bmatrix}
      0 & 1 & 1
  \end{bmatrix}
  -
  \begin{bmatrix}
    0 & 1 & 3
  \end{bmatrix}
  \pause
  \;=\;
  \begin{bmatrix}
    0 & 0 & -2
  \end{bmatrix}
\]

\pause
\[
  \begin{bmatrix}
    1 & 2 & -1 \\
    0 & 1 & 3  \\
    0 & 0 & -2
  \end{bmatrix}
\]
\end{resolution}

%------------------------------------------------------------------------------%
\begin{resolution}{Solução}
Portanto, uma forma escalonada de $A$ é
\(
  \begin{bmatrix}
    1 & 2 & -1 \\
    0 & 1 & 3  \\
    0 & 0 & -2
  \end{bmatrix}
\)
\end{resolution}

%------------------------------------------------------------------------------%
\begin{resolution}{Notação Resumida}
\[
\begin{NiceArray}{rrrl}[columns-width=3ex]
    1 &  2 & -1  & \\
    2 &  5 &  1  & \\
   -1 & -1 &  2  &
  \CodeAfter
    \SubMatrix[{1-1}{3-3}]
  \end{NiceArray}
\]

\pause
\[
  \begin{NiceArray}{rrrl}[columns-width=3ex]
    \hpm 1 & 2 & -1 & \\
         0 & 1 &  3 & \quad L_2 \leftarrow L_2-2 L_1 \\
         0 & 1 &  1 & \quad L_3 \leftarrow L_3+L_1
  \CodeAfter
    \SubMatrix[{1-1}{3-3}]
  \end{NiceArray}
\]

\pause
\[
  \begin{NiceArray}{rrrl}[columns-width=3ex]
    \hpm 1 & 2 & -1 & \\
         0 & 1 &  3 & \\
         0 & 0 & -2 & \quad L_3\leftarrow L_3-L_2
  \CodeAfter
    \SubMatrix[{1-1}{3-3}]
  \end{NiceArray}
\]
\end{resolution}

%------------------------------------------------------------------------------%
